\section{讨论}
\subsection{方法意义}
本文工作的核心意义，在于把研究重心从“进一步增强编码器主干”转向“如何让已有强表示产生更可靠的最终输出”。在临床长文本场景中，这一转向具有现实必要性。许多任务的瓶颈并不只来自特征不足，还来自样本噪声强、标签边界模糊以及模型对困难样本缺乏自我识别能力。因而，当长上下文编码器已经能够提供较稳定的语义表示时，围绕输出层开展可信建模并不是边缘问题，而是决定模型能否进一步走向可解释和可控的重要步骤。

从方法结构上看，\method{} 结合了深度表示与结构化统计估计两种思路。前者负责提取复杂语义特征，后者负责在最终映射阶段显式利用样本权重信息。这样的组合使模型既保留了神经网络的表达能力，又在输出层引入了更明确的归纳偏置。相较于单纯加深 MLP 头，本文方法更强调“为什么某些样本应当被更谨慎地对待”，因此其改进方向更贴近临床应用中的风险控制需求。

\subsection{对参考结果的解释}
从当前参考性结果看，完整 \method{} 在 weighted F1、micro F1、macro AUROC、MSE 和 MAE 上均优于线性头、MLP 头以及阶段 A 版本。这种多指标上的一致改善，说明方法收益并非来自某一特定阈值调节，而更可能反映了输出映射质量的整体提升。尤其是阶段 B 相对于阶段 A 的进一步增益，表明方差分支学习到的样本差异信息能够通过 WLS 精修转化为可观测的性能改进，而不是停留在辅助可视化层面。

消融结果进一步支持了这一判断。去掉方差分支后，模型失去了对样本级噪声的显式刻画；取消 WLS 精修后，方差虽然仍可被预测，但无法再进入最终参数重估过程；将稳定化损失替换为普通异方差 NLL 后，模型在训练可控性和最终指标上均出现回落。由此可以看出，本文方法的效果并非来源于孤立模块，而是来自“异方差建模、稳定训练与结构化精修”之间的协同。

\subsection{局限性分析}
尽管本文方法展现出一定潜力，但其适用边界仍需明确。首先，当前框架建立在冻结特征表示之上，如果编码器对目标任务的表征能力本身不足，那么输出层再精细也难以突破由上游表示决定的性能上限。其次，WLS 精修依赖特征矩阵的条件数和正则化设置，在样本量有限、特征高度相关或超参数选择不当时，解析求解仍可能受到数值稳定性限制。再次，本文建模的预测方差主要反映输入相关噪声，它可以作为风险信号，但尚不能直接替代严格校准后的概率解释。

此外，本文当前实验部分仍以阶段性结果为主，相关趋势尚需在完整的数据整理、统一脚本和多随机种子设置下进一步验证。因此，现阶段更稳妥的表述是：本文提供了一条可行的方法路径和一套自洽的分析框架，而不是对临床可信预测问题的最终结论。

\subsection{对后续工作的启示}
本文的一个直接启示是，可信预测不必从高成本的大模型重训练开始，也可以先在输出层建立噪声感知和结构化修正能力。这一路线对于资源受限但希望形成明确方法贡献的研究工作尤其有参考价值。沿着当前框架，后续可以自然扩展到三个方向：一是比较不同池化策略、损失函数和正则参数对方差分布与主任务性能的影响；二是引入 Bayesian Last Layer、conformal prediction 等机制，使输出从“可给出方差”进一步走向“可校准、可拒识、可审计”；三是在保持成本可控的条件下探索局部参数高效微调，从而研究表示提升与输出层可信建模之间的协同关系。

\subsection{应用与伦理边界}
临床文本研究天然涉及隐私保护、误用风险和责任边界。即便模型能够给出不确定性信号，也不能据此将其视为可直接替代临床判断的自动决策系统。更合理的定位是，将其作为风险提示和人工复核辅助工具，用于帮助医生识别高风险样本、调整审阅优先级或发现可能需要再次核查的病例。只有在数据合规、评估充分和工作流边界清晰的前提下，可信输出层方法才可能真正服务于临床应用。
